% !TEX root = ../main.tex
\chapter{서론과 연구 문제}
\label{ch:introduction}

디파이(DeFi, 탈중앙 금융)는 누구나 볼 수 있는 블록체인 위에서 돈(자본)이 쓰일 곳을 정하는 방식을 크게 바꾸었어요. 이제는 예전의 은행 같은 중간 기관 없이도 돈을 빌려주고, 빌리고, 빌린 돈을 보태 거래할 수 있어요 \parencite{schar2021defi}. 그런데 신용 거래 비슷한 일은 아직도 대부분 담보를 많이 맡겨야만 할 수 있어요. 디파이 서비스(프로토콜)는 블록체인 밖의 신분 확인 제도나 이미 쌓인 신용 기록을 쓸 수 없기 때문이에요 \parencite{harvey2021defi}. 진짜 이름 대신 주소로만 활동하는 세상(가명성)에서는, 서로 모르는 사람들 사이에 돈을 어떻게 잘 맡길지가 가장 큰 문제예요.

이 논문은 이런 것을 물어요. 걷기(화살표를 따라 점에서 점으로 옮겨 다니는 계산) 하나가 ERC-20(토큰 공통 규칙) 허락(allowance, 토큰을 얼마까지 대신 써도 된다는 허락) 화살표(엣지)를 송금(transfer)과 함께 따라가면, 그 허락 화살표는 온체인(블록체인 위의) 지갑 평판을 재는 페이지랭크(PageRank)에 무엇을 보탤까요? 허락 화살표는 \texttt{approve}(허락하기)와 \texttt{permit}(서명 허락)이 남기는 \texttt{Approval} 이벤트(계약이 남기는 기록)에서 나와요. 이 논문의 방법은 결합 페이지랭크(C-PR, 두 그래프를 함께 걷는 점수)예요. C-PR은 같은 주소들의 허락 화살표와 송금 화살표를 하나의 사슬(끊기지 않고 이어지는 한 번의 걷기) 안에서 함께 걸어요. 그리고 주소마다 미리 정해 둔 섞는 비율 $\lambda$로 두 화살표를 섞어요. 한 층(그래프 하나)만 걷는 두 걷기가 C-PR의 비교 대상이에요. 하나는 엔도스랭크(EndorseRank)예요. 엔도스랭크는 허락 화살표만으로 계산한 페이지랭크이고, 화살표 무게(가중치)는 \textcite{do2023awp}의 AWP(적응형 가중 페이지랭크)에서 가져와요. AWP는 송금을 바탕으로 한 방법이에요. 다른 하나는 송금 위에서 계산한 AWP 그 자체예요. 이 점수들은 모두 아비트럼 원(Arbitrum One)의 같은 스마트 계약(블록체인 위에서 저절로 실행되는 프로그램) 계정들에 대해 계산해요. 그 곁에서, 걷기들이 매끄럽게 다듬는(평활하는) 원시 차수(화살표를 그냥 센 수)도 함께 계산해요. 그리고 AWP를 검사할 때 쓴 순위 상관 방법(두 순위가 얼마나 비슷한지 재는 방법)으로 평가해요 \parencite{do2023awp}. 이 연구에서는 그 방법에 짝지은 대비(같은 대상에서 두 점수를 맞대어 보는 비교), 원시 차수 기준선(비교 기준), 시간 홀드아웃(나중 기간을 떼어 두고 맞혀 보는 시험)을 더했어요. 이 논문에서 온체인 평판 지표는, 오프체인(블록체인 밖의) 신원 정보를 하나도 쓰지 않고 온체인 활동과 주고받은 일에서 끌어낸 값이에요 \parencite{do2023hodgerank}.

\dissertationSubheading[sec:background]{배경: 디파이, 담보, 가명성}

누구나 볼 수 있는 블록체인은 중앙 기관 없이도 프로그램으로 거래를 마무리(결제)할 수 있게 해 줘요 \parencite{nakamoto2008bitcoin}. 스마트 계약은 약속의 조건을 코드로 적어 둘 수 있어요 \parencite{szabo1997formalizing}. 그리고 이런 계약을 쓰는 디파이 플랫폼은 돈을 맡아 주는 기관(수탁 기관)을 거치지 않고 블록체인 위에서 거래를 마무리해요 \parencite{buterin2014ethereum}. 그다음 서비스들은 이 계약들을 쌓아 올려, 빌려주기(대출), 자동 마켓 메이킹(프로그램이 자동으로 사고팔아 주는 일), 무기한 선물 거래를 하는 탈중앙 시장을 만들어요 \parencite{werner2022sok}. 아비트럼 원과 다른 레이어 2(L2) 네트워크에서는 수수료가 충분히 낮아서, 보통 사용자도 자주 거래해요. 그리고 모든 송금, 허락, 서비스 이벤트를 지갑 하나하나마다 살펴볼 수 있어요.

이렇게 다 보이는데도, 디파이의 신용 시장은 본래부터 과담보(빌리는 돈보다 담보를 더 맡기는) 상태이고, 지금도 그대로예요. 대부분의 서비스는 빌리는 사람이나 트레이더(거래하는 사람)에게 포지션(거래로 걸어 둔 몫)의 값보다 더 많은 담보를 맡기게 해요. 거래 상대가 가명의 계정이고, 그 계정의 주인이 누구인지 확인되지 않았기 때문이에요 \parencite{gudgeon2020crisis,harvey2021defi}. 가명성은 누구나 볼 수 있는 장부의 기본 성질이에요. 누구나 거의 돈을 들이지 않고 주소를 만들 수 있어요. 그리고 한 주체(한 사람이나 한 단체)가 지갑을 둘 이상 가질 수 있어요 \parencite{douceur2002sybil}. 보통의 신용 등급은 개인의 지난 기록에 기대는데, 그런 기록은 이런 환경에 잘 옮겨지지 않아요 \parencite{packin2024credit}. 그래서 서비스들은 한 사람 한 사람에게 맞춘 평판 대신, 쌓아 둔 담보, 청산(손실이 너무 커져 거래가 강제로 끝나는 일) 장치, 서비스 규칙으로 위험을 다스려요.

담보를 맡기게 하는 것은 서비스를 지켜 주지만, 돈(자본)을 알뜰하게 쓰지 못하게 해요. 트레이더와 빌리는 사람은 다른 곳에 투자할 수도 있었던 돈을 묶어 둬야 해요. 그리고 전통 금융은 신용 기록이 있는 사람에게 담보를 빌리는 돈보다 적게 맡겨도 되는 상품을 내주지만, 이런 플랫폼은 그런 상품을 하나도 내놓지 않아요. 그래서 연구자들은 오프체인 신용 기록 대신 쓸 것으로 온체인 신호에 눈을 돌렸어요. 그런 연구로는 송금 그래프(점과 화살표로 이은 그림)의 중심성(그래프에서 얼마나 가운데에 있는지) \parencite{do2023hodgerank}, 네트워크를 따라 위험이 퍼지는 것(위험 전파)을 보는 방법 \parencite{lin2025riskprop}, 전망 이론(사람이 얻는 것과 잃는 것을 어떻게 느끼는지 다루는 이론)과 블록체인에 바탕을 둔 신용 점수 매기기 \parencite{hassija2020lending}, 대출 서비스를 대상으로 한 기계 학습 처리 과정(파이프라인) \parencite{kotb2023credit}이 있어요. 이 연구 흐름들은 하나하나 모두, 공개된 장부 자료에 사용자가 누구인지 밝혀내지 않고도 위험 평가에 쓸 수 있는 정보가 들어 있다고 가정해요.

그래프 순위 알고리즘은 이 문제에 잘 맞아 보여요. 페이지랭크와 그 변형들은 되풀이되는 보증(믿고 맡김)으로 점(노드)에 순위를 매겨요. 그래서 한 점의 점수는 다른 점들에게서 받는 보증에 달려 있어요 \parencite{page1999pagerank}. 예전에 나온 웹 신뢰 알고리즘인 아이겐트러스트(EigenTrust) \parencite{kamvar2003eigentrust}와 트러스트랭크(TrustRank) \parencite{gyongyi2004trustrank}는 비슷한 생각을 컴퓨터끼리 직접 주고받는 망(피어 투 피어)의 신뢰와 웹 스팸 문제에 써요. \textcite{do2023hodgerank}에서는 송금 그래프에 페이지랭크와 호지랭크(HodgeRank)를 함께 써서 이더리움(Ethereum) 사용자가 얼마나 믿을 만한지 판단해요. 그리고 그 결과를 일종의 사회적 신용이라고 불러요. AWP도 송금 그래프로 순위를 매기는 또 다른 알고리즘이에요 \parencite{do2023awp}. AWP는 송금마다 무게를 매기는데, 금액에 대한 상한이 있는 함수와 송금의 나이(얼마나 오래되었는지)에 대한 로지스틱 함수(오래된 일일수록 가볍게 세는 함수)를 써요. 그리고 최근에 송금을 보낸 주소에서 걷기를 다시 출발(재시작)해요. AWP를 만든 사람들은 AWP 점수가 들어오는 송금의 수와 금액과 함께 오르내린다는 것을 보여서 AWP를 검증했어요. 이때 들어오는 송금의 수와 금액은 경제적 중요도의 대리 지표(대신 재는 값)로 보았어요. 이 연구들 덕분에 송금 그래프 평판은 믿을 만한 기준선이 되었어요. 하지만 이 연구들은 쓰는 권한을 분명하게 내주는 일을 빼놓아요. 이 일은 송금과 따로 일어나는 온체인 행동이고, ERC-20 허락으로 기록돼요 \parencite{vogelsteller2015erc20,lundfall2020permit}.

\dissertationSubheading[sec:motivation]{온체인 평판이 필요한 까닭}

레버리지(빌린 돈을 보태 거래하는 것)를 쓰는 디파이에서 온체인 평판 연구가 필요한 주된 까닭은 네 가지예요.
\begin{itemize}
\item 자본 효율(돈을 알뜰하게 쓰는 정도): 서비스들은 순서를 나타내는 순위 신호를 찾고 있어요. 이런 신호는 등급을 나눈 서비스(예를 들어 한도를 더 자주 올려 주기, 수수료 낮춰 주기, 감시할 때 더 높은 우선순위 주기)를 뒷받침할 수 있어요. 어쩌면 앞으로 할 통제된 실험(조건을 맞춰 두고 하는 실험)에서는 담보 조건을 서로 다르게 정하는 데까지 쓸 수 있을지도 몰라요. 이때도 거래의 마무리는 온체인에 그대로 남아요. 온체인에서는 계약이 담보 비율과 청산 조건을 지키게 해요 \parencite{schar2021defi}.
\item 검사할 수 있음: 오프체인 신용 점수는 속을 들여다볼 수 없어요. 하지만 공개된 이벤트 자료로 계산한 그래프 점수는 다시 계산하고 검사할 수 있어요. 이것은 속을 알 수 없는(블랙박스) 탈중앙 신용 점수 매기기에 대한 걱정을 어느 정도 덜어 주는 데 도움이 될 거예요 \parencite{packin2024credit}.
\item 운영 비용: 예를 들어 레버리지를 바꾸기 전마다 새로 계산하는 것처럼, 자주 새로 고치는 점수는 다시 계산하는 데 비용이 적게 들어야 해요.
\item 표준화: 다른 시장에서는 거래 상대를 가려내려고 심사 비율(상대를 걸러 내는 숫자 기준)을 써요. 하지만 가명의 디파이에는 지갑 단위에서 이 일을 할, 순서를 나타내는 표준 언어가 없어요. 이 논문에서 살펴본 점수들이 어디에서 그 역할을 할 수 있고 어디에서는 할 수 없는지는 \ref{ch:discussion}장에서 이야기해요.
\end{itemize}

송금에 바탕을 둔 AWP는 공개되어 있고 검사할 수도 있어요. 하지만 AWP의 그래프는 돈이 오갈 때마다 커져요 \parencite{do2023awp}. 허락 그래프는 짜임새가 다른 또 하나의 길이에요. 사용자가 ERC-20 토큰을 허락할 때마다, 계약은 스펜더(spender, 허락을 받아 대신 쓰는 쪽)가 그 토큰을 얼마까지 옮길 수 있는지 기록해요 \parencite{vogelsteller2015erc20,lundfall2020permit}. 그리고 새 허락은 예전 허락을 덮어써요. 그래서 주인(허락을 해 주는 쪽)--스펜더 쌍과 토큰마다, 마지막 허락(가장 최근 허락) 하나가 그 쌍의 지난 기록을 대신해요. 허락은 송금보다 드물기도 해요. 그래서 같은 계약들에서 보증 그래프 $G_{\text{approve}}$의 화살표 수는 송금 그래프 $G_{\text{transfer}}$의 절반도 안 돼요. 페이지랭크를 푸는 계산은 반복(한 바퀴)마다 드는 비용이 화살표 수 $|E|$에 정비례해서 커지는데, 이 계산에 걸리는 시간도 더 짧아요(표~\ref{tab:benchmark-runtime}). 두 경우 모두 같은 풀이 프로그램(solver)이 돌아가요. 그래서 이 속도 향상(몇 배 빠른지)은 더 듬성듬성한(화살표가 적은) 그래프와, 그 그래프에서 걷기에 필요한 반복 횟수에서 오는 것이고, 알고리즘이 나아진 것은 아니에요. 주장~C5는 이것을 그래프를 나타내는 방식이 가진 운영상의 성질로 다뤄요.

효율 말고도, 허락은 뜻도 달라요. 송금은 다른 곳으로 가는 길에 거쳐 가는 것일 수도 있고, 마켓 메이킹일 수도 있고, 그냥 가만히 받은 것일 수도 있어요. 반면 허락은 돈을 써도 된다는 권한을 일부러 내주는 일이에요. 페이지랭크는 화살표를 인용(다른 글을 가리켜 추천하는 것)처럼 읽어요 \parencite{page1999pagerank}. 허락을 보증 화살표로 다루면, 이 읽는 법을 서비스들이 이미 기대고 있는 믿음 관계로 옮겨 올 수 있어요. 이렇게 하면 순위가 달라지는지, 그 순위가 허락 점검과 송금 점검에 맞는 정도가 달라지는지, 그리고 두 종류의 화살표를 모두 걷는 걷기 하나가 어느 한 종류만 걷는 걷기보다 계약에 순위를 더 잘 매기는지를 살펴봐요. 살펴보는 대상은 아비트럼 원의 매칭 코호트(모든 점수가 똑같이 점수를 매기는 계약 묶음)에 든 계약 계정 $n=6{,}483$개예요.\footnote{매칭되었다는 말은 모든 방법이 같은 계약들에 점수를 매긴다는 뜻이에요. 그 계약들은 2025년 7월 1일부터 2026년 6월 30일까지 WETH, USDC, USDT0, ARB, WBTC 가운데 어느 토큰에서든 서로 다른 주인 셋 이상에게서 허락을 받은 계약이에요. 이 목록은 허락 로그(기록)만으로 뽑았어요(\ref{sec:data-sources}절).}

\dissertationSubheading[sec:research-problem]{연구 문제}

AWP는 송금 기록에서 무게를 단 합과 로지스틱 감쇠(오래된 일일수록 가볍게 세는 무게)를 써서 믿음을 짐작해요 \parencite{do2023awp}. AWP는 (i) 송금 기록 전체를 모아 계산해야 하고, (ii) 값진 지갑을 들어오는 송금을 많이 받는 지갑으로 보며, (iii) 쓰는 권한을 내주는 일에는 아무 역할도 주지 않아요. 이 논문의 허락 쪽 비교 대상인 엔도스랭크는 AWP의 화살표 무게와 풀이 프로그램을 그대로 써요. 하지만 방향이 있는 그래프는 주인--스펜더 쌍과 토큰마다의 마지막 허락으로 만들어요. 그렇게 만든 그래프는 더 듬성듬성하고, 다른 신호를 담고, 각 관계에서 다른 쪽에 순위를 매겨요. 허락은 주인에서 스펜더 쪽을 가리켜요. 그래서 엔도스랭크는 라우터나 금고처럼 허락을 받는 스펜더에 점수를 매기고, 허락을 하나도 받지 않은 주인에게는 모두 같은 값을 줘요. 이렇게 해서 두 걷기는 저마다 한 가지 행동만 따라가요. AWP는 한 주소에 닿는 돈(지급)을 따라가고, 엔도스랭크는 그 주소가 받은 쓸 권리를 따라가요.

두 종류의 화살표는 같은 주소들 위에 있어요. 그래서 걷기 하나가 두 화살표를 모두 따라갈 수 있고, 미리 정해 둔 섞는 비율 $\lambda$로 점마다 둘을 섞을 수 있어요. 이 걷기가 결합 페이지랭크(C-PR)이고, 이 논문의 방법이에요. $\lambda=1$과 $\lambda=0$에서는 C-PR이 한 층만 걸어요. 두 종류의 화살표는 같은 계약들에 대해 서로 다른 행동을 기록해요. 그래서 이렇게 섞은 것을 시험해 볼 가치가 있어요. 이미 송금 걷기를 돌리고 있는 운영자는 점수를 하나 더 돌리는 대신, 그 걷기에 허락 화살표를 더할 수 있어요. 주된 연구 질문(RQ)은 이거예요. ERC-20 허락 화살표와 송금 화살표를 모두 걷는 페이지랭크 걷기 하나는, 아비트럼 원의 계약 계정에 순위를 매기는 데 무엇을 보탤까요? 어느 한 종류의 화살표만 걷는 걷기와 비교하면, 그리고 그 걷기들이 매끄럽게 다듬는 원시 차수와 비교하면 어떨까요? RQ1--RQ5는 이 질문을 잘게 나눈 연구 질문이에요(\ref{sec:research-questions}절). 분석 계획은 판정 규칙 하나를 미리 정해 두었고, 그 규칙의 대상은 C-PR이에요(\ref{sec:holdout-protocol}절).

점수가 자기 화살표의 통계와 잘 맞거나, 같은 화살표가 앞으로 더 생길 것을 예측(미리 맞히기)한다고 해도, 그것만으로는 그 점수가 알맹이 있는 뜻에서 평판을 잰다고 아직 말할 수 없어요. 그래서 두 번째 홀드아웃 기간과 함께 등록(무엇을 어떻게 시험할지 미리 정해 기록해 두는 일)한 비교 하나는, 두 종류의 화살표 어느 것으로도 만들지 않은 결과도 써요. 이 비교는 연구 질문을 새로 더하지 않아요. 그 결과는 동결(점수를 고정하는 일) 뒤에 관측했어요. 바로 GMX~V2(무기한 선물 거래 서비스) 계정 가운데 계약인 계정들이 그 뒤에 청산을 피했는지예요. 이 비교의 표본은 작아요(\ref{sub:neutral-label-protocol}절). 이 비교는 고친 점수(등록한 뒤에 바꾼 방식으로 다시 매긴 점수)에 재분석으로 적용해요.

점수들은 표~\ref{tab:primary-methods-preview}에 있어요. 표~\ref{tab:hrq-claim-map}에는 가설(H1--H3), 연구 질문(RQ1--RQ5), 그리고 \ref{ch:results}장에서 나온 실험 결과 주장(C1--C6)을 정리했어요.

\begin{table}[htbp]
\centering
\caption{이 논문에서 평가한 점수들(매칭 계약 코호트, $n=6{,}483$).}
\label{tab:primary-methods-preview}
\fitwidth{%
\begin{tabular}{p{0.18\linewidth}p{0.28\linewidth}p{0.48\linewidth}}
\toprule
점수 & 입력 그래프 & 역할 \\
\midrule
C-PR & 두 층 모두, 달러(USD) 금액(허락은 하나하나 상한으로 자름), 점마다 섞는 비율 $\lambda$로 섞음 & 이 논문의 방법이고, 규칙~A와~B의 대상이에요. 한 층만 걷는 C-PR의 층들, 곧 C-PR ($\lambda=1$)과 C-PR ($\lambda=0$)은 금액으로 무게를 달아요. 그래서 이 둘은 EndorseRank, AWP와 같지 않아요. 이 점수의 빼 보기 실험(한 부분을 빼고 다시 해 보는 실험)인 S-PR(보증 씨앗 페이지랭크)은 송금 걷기를 허락 걷기에서 다시 출발시켜요 \\
EndorseRank & 다섯 토큰의 마지막 ERC-20 허락, 달러 금액($G_{\text{approve}}$), AWP처럼 무게를 매김 & 허락으로 맡김; 허락 위의 한 층 비교 대상이에요. 허락을 해 주는 주인에서 다시 출발하는 변형도 옆에 함께 보고해요 \\
AWP & 다섯 토큰의 송금, 달러 금액($G_{\text{transfer}}$) & 송금 흐름; 송금 위의 한 층 비교 대상이고, \textcite{do2023awp}의 발표된 송금 그래프 방법이에요(\ref{sub:awp-graph}절) \\
원시 차수 & $G_{\text{approve}}$와 $G_{\text{transfer}}$, 점마다 셈 & 받은 허락 수(허락 입차수)와 송금 입차수(송금을 보내온 서로 다른 주소 수). 페이지랭크들이 매끄럽게 다듬는 그 점 둘레의 수예요. 같은 기간에서는 대리 지표, 기간 밖에서는 기준선으로 써요 \\
\bottomrule
\end{tabular}
}
\end{table}

주요 용어의 뜻은 다음과 같아요.

\begin{itemize}
\item \emph{디파이(DeFi, 탈중앙 금융):} 빌려주기와 빌리기 같은 금융 서비스예요. 누구나 볼 수 있는 블록체인 위에 스마트 계약으로 만들어져 있어요. 그리고 대개 돈을 맡아 두거나 거래를 마무리하는 일을 책임지는 중앙 운영자 없이 돌아가요 \parencite{schar2021defi}.
\item \emph{스마트 계약:} 블록체인 위에서 스스로 실행되는 프로그램 코드예요. 믿고 맡길 중간 기관 없이 약속을 실행해요 \parencite{szabo1997formalizing}.
\item \emph{지갑 / 주소:} 돈을 담고 스마트 계약과 주고받는 블록체인 계정이에요. 이 계정은 자기 주소로 알아볼 수 있어요 \parencite{buterin2014ethereum}. 지갑은 둘 가운데 하나예요. 하나는 외부 소유 계정(EOA)으로, 개인 키(주인만 아는 비밀 값)로 움직이는 사용자의 지갑이에요. 다른 하나는 자기 코드에 따라 움직이는 계약 계정이에요 \parencite{antonopoulos2018mastering}. 여기서 순위를 매기는 지갑은 계약 계정이에요. EOA는 주인, 보낸 쪽(송신자), 받는 쪽(수신자)으로만 들어와요. 오프체인 신원은 미리 가정하지도 않고 알아내지도 않아요.
\item \emph{온체인 평판 점수:} 온체인 자료로 계산한 수예요. 가명성 아래에서 금융 활동을 할 때, 한 지갑이 상대적으로 어디쯤에 있는지 나타내요 \parencite{do2023hodgerank}. 여기서 점수를 받는 주소는 계약이에요. 오프체인에서 가져오는 입력은 다섯 토큰의 목록과 하루 달러 가격뿐이에요. 점수는 순서를 나타내는 값이고, 순위를 매기려고 써요. 점수 크기 자체는 신용 한도를 정하는 데 아무 뜻이 없어요.
\item \emph{페이지랭크:} 감쇠 계수(화살표를 따라갈 확률)를 써서 방향이 있는 그래프 전체로 영향력을 퍼뜨리는 그래프 순위 알고리즘이에요 \parencite{page1999pagerank}. 적응형 가중 페이지랭크(AWP)는 송금 그래프 위의 페이지랭크예요. 화살표마다 송금의 금액과 오래된 정도로 무게를 매기고, 최근에 활발했던 보낸 쪽에서 다시 출발해요 \parencite{do2023awp}. 여기서 금액은 달러로 나타내요. 그리고 금액 무게는 금액이 중앙값일 때 2분의 1이에요. 엔도스랭크는 AWP의 화살표 무게를 마지막 허락 그래프에 쓰고, 고르게 다시 출발(균등 재시작)해요(\ref{sec:graph-construction}절).
\item \emph{결합 페이지랭크(C-PR):} 이 논문의 방법이에요. 허락 층과 송금 층을 함께 걷는 무작위 걷기 하나예요. 각 지갑에서는 그 지갑이 들어 있는 층들에 대해, 두 층의 행(이동 행렬에서 그 지갑의 줄)을 무게 $\lambda$와 $1-\lambda$로 섞어요(\ref{sub:coupled-graph}절). 보증 씨앗 페이지랭크(S-PR)는 송금 걷기를 그대로 두지만, 허락 걷기 결과에서 다시 출발해요. 트러스트랭크가 씨앗 집합(믿을 만하다고 미리 고른 점들)에서 다시 출발하는 것과 같아요 \parencite{gyongyi2004trustrank}.
\item \emph{ERC-20 허락 / Approve / Permit:} 스펜더가 주인 대신 주인의 토큰을 얼마나 쓸 수 있는지 기록하는 계약 함수예요 \parencite{vogelsteller2015erc20,lundfall2020permit}. \texttt{Approve}는 보통의 온체인 호출이에요. EIP-2612 \texttt{permit}은 서명으로 같은 상태 변화를 만들어요. 엔도스랭크는 주인--스펜더 쌍과 토큰마다, 0이 아닌 마지막 허락을 써요.
\item \emph{분야 일치도:} 평판 순위가 눈으로 볼 수 있는 대리 지표의 순위와 얼마나 맞는지 나타내는 정도예요. 스피어먼의 로($\rho$) \parencite{spearman1904proof}와 켄달의 타우($\tau$, 두 줄 세우기가 얼마나 같은 순서인지 나타내는 수) \parencite{kendall1938rank}로 재요. 그리고 구성 타당도(점수가 재려는 것을 정말 재는지)에 대한 증거로 읽어요 \parencite{cronbach1955construct}. 점수 자신의 화살표로 만든 대리 지표와 견줄 때는, 이것은 의도한 구성개념(점수가 재려고 하는 생각) 점검일 뿐이에요.
\end{itemize}

\dissertationSubheading[sec:supplementary-extension]{시간 홀드아웃}

같은 기간에서 잰 허락 쪽 켄달 $\tau$(표~\ref{tab:alignment-hybrid})는 의도한 구성개념 점검이에요. $G_{\text{approve}}$ 위의 페이지랭크는 스펜더가 받은 허락 수와 잘 맞아야 해요. 페이지랭크는 그 수를 그래프 전체에 걸쳐 매끄럽게 다듬기 때문이에요. 점수가 같은 기간의 그래프 너머의 정보를 전하는지 시험하려고, 점수를 $t_1$(2026년 3월 31일, 23:59:59~UTC)에 동결해요. 그리고 2026년 4월--6월이 되어야 얻을 수 있는 두 라벨(나중에 채점에 쓰는 정답 값)과 견주어요. 하나는 새 주인--스펜더 허락 쌍이에요(보증 구성개념). 다른 하나는 새 송금자(새로 송금을 보내온 주소)예요(흐름 구성개념). 이 첫 번째 기간인 W0에서, 홀드아웃 코호트($n=7{,}914$)는 뽑아낸 허락 자료에서 $t_1$ 시점에 다섯 토큰 가운데 하나의, 0보다 큰 마지막 허락을 적어도 하나 가진 모든 스펜더 계약이에요. 이 마지막 허락은 $t_1$ 전 열두 달 안에 정해진 것이에요. 또 각 페이지랭크를 그것이 $t_1$에 매끄럽게 다듬는 원시 차수와도 견주어요.

두 번째 기간인 W1은 2026년 10월 9일 16:49~UTC에 등록했어요. W0 결과를 안 뒤였고, 이 연구가 2026년 6월 30일 뒤 날짜의 로그를 하나도 뽑기 전이었어요. W1의 점수는 2026년 6월 30일에 동결하고, 라벨은 2026년 7월부터 9월까지 세요. W1은 그 동결일의 스펜더 계약 $n=7{,}960$개에서 두 스펜더 라벨을 되풀이해요. 그리고 계약인 GMX~V2 트레이더($n=70$)에 대한 라벨 하나를 더해요. 바로 그 트레이더들이 포지션을 닫은 것 가운데 청산으로 끝나지 않은 것의 비율이에요(\ref{sub:fresh-holdout}절).

2026년 10월 9일의 분석 계획과 W1의 등록이 등록 분석(미리 등록한 그대로 한 분석)을 정해요. 등록 분석에서는 모든 토큰의 금액을 원래 단위(소수점 조정을 하지 않은 가장 작은 단위)로 나타내요. 등록 분석의 결과를 안 뒤에 수정 분석(등록 뒤에 금액을 달러로 바꿔 다시 한 분석)을 정했어요. 수정 분석에서는 널리 거래되는 다섯 토큰의 금액을 달러로 나타내고, 각 동결일 전 열두 달의 자료로 점수를 매겨요(\ref{sec:amount-revision}절). 주된 분석은 수정 분석이에요. 등록 분석의 규칙에 따라 내린 수정 분석의 판정은 재분석이지, 등록한 시험이 아니에요. 그리고 등록 분석은 R1만 빼고 같은 판정에 이르렀어요(부록~\ref{sec:map-registered}).

\begin{table}[htbp]
\centering
\caption{가설, 연구 질문, 자료로 확인한 주장의 요약.}
\label{tab:hrq-claim-map}
\fitwidth{%
\begin{tabular}{p{0.12\linewidth}p{0.42\linewidth}p{0.40\linewidth}}
\toprule
항목 & 다루는 것 & 보고한 곳 \\
\midrule
H1, H2 / RQ1 & 허락 화살표만 걷는 EndorseRank와 송금 화살표만 걷는 AWP는 같은 계약들을 서로 다른 순서로 세우나요? 그리고 둘은 저마다 자기 화살표 종류의 점검과 가장 잘 맞나요? & \S\ref{sec:same-window-alignment}과~\S\ref{sec:rank-divergence}; 표~\ref{tab:alignment-hybrid}과~\ref{tab:tie-sensitivity}; 주장 C4 \\
H3 / RQ2 & $n$이 같을 때 각 걷기를 계산하는 데 얼마나 드나요? 그리고 두 화살표 집합을 모두 걸으면 얼마나 더 드나요? & 표~\ref{tab:benchmark-runtime}과~\ref{tab:benchmark-scaling}; 주장 C5 \\
RQ3 & $t_1$의 점수가 $t_2$의 새 허락과 새 송금자를 예측하나요? 그리고 결합 걷기를 포함해, 같은 화살표 종류의 원시 차수보다 더 잘하는 걷기가 있나요? & \S\ref{sec:holdout-results}과~\S\ref{sub:cpr-contrasts}; 표~\ref{tab:holdout-spenders}, \ref{tab:holdout-diff}, \ref{tab:holdout-receiving}, \ref{tab:holdout-traders}, \ref{tab:fresh-holdout}, \ref{tab:fresh-posthoc}과~\ref{tab:cpr-contrasts}; 주장 C2와 C3 \\
RQ4 & C-PR은 같은 기간에서도 기간 밖에서도, 자기 층 가운데 하나만 걷는 어느 것보다 계약에 순위를 더 잘 매기나요(규칙 A)? & \S\ref{sec:coupled-results}과~\S\ref{sub:cpr-rule-a}; 표~\ref{tab:alignment-hybrid}, \ref{tab:tau-diff-hybrid}와~\ref{tab:holdout-diff}; 주장 C1 \\
RQ5 & 허락 층을 더하면 새 자료에서 송금 걷기가 나아지나요(규칙 B)? & \S\ref{sub:fresh-results}; 표~\ref{tab:fresh-holdout}와~\ref{tab:fresh-contrasts}; 주장 C1 \\
C1 & 고친 점수에 규칙 A와 B를 적용하면, C-PR은 한 층만 걷는 어느 쪽도 이기지 못하고, 허용 폭 안에서 자기 송금 층보다 나아지지도 않아요(새 허락 쌍에서 나은 점은 그대로 나타나지만, 새 송금자에서, 그리고 아슬아슬하게 청산 없이 닫은 비율에서 비열등성(정한 폭보다 더 나쁘지 않음)을 채우지 못해요). AWP와 견주면 뚜렷하게 실패해요. S-PR에서 허락 화살표는 주로 걷기를 직접 이끌 때 도움이 돼요 & 표~\ref{tab:tau-diff-hybrid}, \ref{tab:holdout-diff}과~\ref{tab:fresh-contrasts} \\
C2 & 규칙과 별도로 보면, C-PR은 두 기간 모두에서 두 스펜더 라벨을 EndorseRank보다 더 잘 줄 세워요. 또 새 허락 쌍을 AWP보다 더 잘 줄 세우는데, 이것은 대부분 AWP가 0점을 매기는 스펜더들 덕분이에요. 반면 새 송금자는 AWP와 송금 입차수보다 못하게 줄 세워요. 새 허락 쌍은 받은 허락 수보다 못하게 줄 세우고, W1에서는 활동 재시작(활발한 주소에서 더 자주 다시 출발하기)을 쓴 EndorseRank보다도 못하게 줄 세워요(라벨을 본 뒤에 계산한 대비) & 표~\ref{tab:cpr-contrasts}; \S\ref{sub:cpr-contrasts} \\
C3 & 기간 밖에서는 한 층 걷기 어느 것도 자기가 다듬는 원시 차수를 이기지 못해요. 두 기간 모두에서 AWP는 나중의 새 송금자를 송금 입차수보다 못하게 줄 세우고, EndorseRank는 어느 다시 출발 규칙을 쓰든 나중의 새 허락 쌍을 받은 허락 수보다 못하게 줄 세워요(라벨을 본 뒤에 계산한 대비) & 표~\ref{tab:holdout-spenders}, \ref{tab:holdout-diff}, \ref{tab:holdout-receiving}과~\ref{tab:fresh-posthoc} \\
C4 & 같은 기간에서 한 층 걷기마다 자기 화살표 종류의 점검과 잘 맞아요(일부는 만든 방식 때문에 저절로 그래요). 그리고 두 걷기는 코호트를 서로 다른 순서로 세워요 & 표~\ref{tab:alignment-hybrid}과~\ref{tab:tie-sensitivity}; 그림~\ref{fig:rank-scatter} \\
C5 & 각 점수의 실행 시간과 그래프 크기: 속도는 알고리즘이 아니라 그래프 크기와 각 걷기에 필요한 반복 횟수를 따라가요. C-PR과 S-PR은 AWP보다 느려요 & 표~\ref{tab:benchmark-runtime} \\
C6 & GMX~V2 계약 트레이더가 나중에 청산을 피하는 것에서, 등록 기간(이 연구가 자료를 뽑기 전에 등록해 둔 기간)에는 그냥 센 수의 수준에서 허락이 송금보다 앞서요(W0에서는 양수이지만 어느 쪽이 나은지 가려지지 않아요). 페이지랭크 버전들은 어느 쪽이 나은지 가려지지 않아요. W1에서 이익을 낸 비율과 손실이 아닌 비율에서는 허락한 주인 수(서로 다른 주인 수)가 더 낮아요. 표본이 작아요 & 표~\ref{tab:neutral-label-contrasts}와~\ref{tab:neutral-label-recipients} \\
\bottomrule
\end{tabular}
}
\end{table}

\dissertationSubheading[sec:research-objectives]{연구 목표}

이 연구에는 목표(O)가 다섯 개 있어요. RQ1--RQ5마다 하나씩이에요. 각 목표는 그것을 실제로 재는 데 쓴 측정과 함께 적었어요. 또 \ref{ch:discussion}장에서 그 목표를 이루었는지 판단할 때 쓰는 기준과 함께 적었어요.

\begin{enumerate}
\item[O1.] 두 한 층 걷기인 엔도스랭크와 AWP가 같은 코호트의 계약들을 서로 다른 순서로 세우는지 밝혀요. 그리고 같은 기간에서 각 점수가 각 구성개념과 얼마나 잘 맞는지 평가해요(RQ1). 측정: 두 점수 벡터 사이의 켄달의 $\tau$($n=6{,}483$)와 그 95\% 부트스트랩(다시 뽑기) 구간이에요. 또 모든 점수와 구성개념에 대해 구한, 지표 가족(비슷한 지표 묶음) 수준의 켄달의 $\tau$와 95\% 구간이에요. 기준: 두 걷기 사이 일치의 구간은 $0$(서로 상관없는 순서)도, $1$(똑같은 순서)도 품지 않아야 해요. 어떤 점수의 그 가족에 대한 구간이 $0$을 품지 않으면, 그 점수는 그 구성개념과 맞는다고 봐요. 점수가 자기 화살표의 구성개념과 맞는 것은 의도한 구성개념 점검으로 읽어요. 그 대리 지표들이 그래프와 같은 이벤트에서 나오기 때문이에요.
\item[O2.] 같은 계약들에서 각 점수를 계산하는 데 얼마나 드는지 재요. 여기에는 두 그래프를 합친 그래프(합집합) 위의 C-PR도 들어가요. 그리고 뽑은 계약 수에 따라 실행 시간이 어떻게 바뀌는지 재요(RQ2). 측정: 푸는 시간(준비 실행을 한 번 한 뒤, 시간 잰 실행 다섯 번의 평균과 표준편차(SD)), 메모리, 반복 횟수, $|V|$와 $|E|$예요. 모두 한 가지 절차로 재요. 기준: 시간의 비율은 $|E|$의 비율과 함께만 보고해요. 그래야 어떤 차이든 그래프의 크기 때문인지 따져 볼 수 있어요.
\item[O3.] $t_1$에 동결한 점수가, 코호트의 스펜더들 사이에서 $t_2$의 새 허락과 새 송금자를 예측하는지 시험해요. 그리고 결합 페이지랭크를 포함해, 같은 화살표 종류의 원시 차수보다 더 잘하는 페이지랭크가 있는지 시험해요(RQ3). 측정: 동결 때 계산한 모든 점수(두 원시 차수 포함)와 두 라벨 사이의 켄달의 $\tau$와 95\% 구간이에요. 이것을 W0(동결 2026년 3월 31일, 라벨 2026년 4월--6월, 스펜더 계약 $n=7{,}914$개)과 W1(동결 2026년 6월 30일, 라벨 2026년 7월--9월, $n=7{,}960$)에서 구해요. 또 한 층 페이지랭크마다 자기 차수와 견주고, C-PR은 두 차수와 모두 견주는 짝지은 대비도 구해요. 기준: 어떤 점수의 구간이 $0$을 품지 않으면, 그 점수는 라벨을 예측한다고 봐요. 페이지랭크는 짝지은 차이의 구간이 $0$보다 위에 있을 때만 그 차수보다 낫다고 봐요. 차수 대비도, 이 기준도 W0 라벨을 알기 전에 정하지 않았어요.
\item[O4.] 두 종류의 화살표를 가로지르는 걷기 하나가, 자기 층을 하나씩만 걷는 것을 저마다 이기는지 시험해요(RQ4). 측정: $\lambda=0.5$인 C-PR($0.25$와 $0.75$는 민감도(조건을 바꿔 결과가 얼마나 움직이는지)를 보는 값)과 그 빼 보기 실험인 S-PR이에요. 이 둘을 한 층만 걷는 두 걷기와 견주어요. 견줄 때는 홀드아웃 라벨과 같은 기간 라벨에서 $\Delta\tau$의 짝지은 구간을 써요. 기준: 규칙 A예요. 규칙 A는 이 연구의 로그를 하나도 뽑기 전에, 2026년 10월 9일의 분석 계획에서 정했어요. 규칙 A에 따르면 C-PR은 W0에서 새 허락에 대해 허락 걷기보다 못하지 않아야 하고, 새 송금자에 대해 송금 걷기보다 나아야 해요. 그리고 같은 기간에서, 송금 가족과 허락 가족에 대해서는 앞선 층의 구간 안에 있어야 하고, 시빌 안정성(가짜 주소로 부풀리기 어렵게 보정한 활동 지표)에 대해서는 두 층보다 모두 위에 있어야 해요. 이 규칙은 고친 점수에 재분석으로 적용해요.
\item[O5.] 송금을 걷는 걷기에 허락 층을 더하면 새 자료에서 그 걷기가 나아지는지 시험해요(RQ5). 측정: W1에서 C-PR을 자기 송금 층만 걷는 것과 견주는 짝지은 대비예요. 새 허락 쌍, 새 송금자, 트레이더의 청산 없이 닫은 비율에서 견주어요. 기준: 규칙 B예요. 규칙 B는 2026년 10월 9일에 등록했어요. W0 결과를 안 뒤였고, 이 연구가 W1 로그를 뽑기 전이었어요. 규칙 B에 따르면 C-PR은 W1에서 새 허락 쌍에 대해 송금 걷기를 이겨야 해요. 그리고 새 송금자나 트레이더의 청산 없이 닫은 비율에서 송금 걷기보다 0.01을 넘게 나쁘지 않아야 해요. 이 규칙은 고친 점수에 재분석으로 적용해요. 등록한 민감도 분석은 송금 걷기 자리에 AWP를 넣고 이 규칙을 되풀이해요. 이 분석은 규칙을 바꾸지 않아요. C-PR을 엔도스랭크, AWP, 다른 점수와 견준 대비는 두 규칙 옆에 함께 보고하지만, 어느 규칙도 판정하지 않아요.
\end{enumerate}

\dissertationSubheading[sec:research-questions]{연구 질문}

첫 두 연구 질문에는 가설(H) 세 개가 따라와요. H1: 같은 계약 묶음에서 엔도스랭크와 AWP는 서로 다른 순서를 만들어요. H2: 각 점수는 자기가 걷는 화살표로 만든 점검과 잘 맞아요. H3: $n$의 값이 같을 때, 점수의 계산 시간은 그 그래프의 크기와 함께 커져요. 그래서 가장 듬성듬성한 그래프 위의 엔도스랭크가 가장 빨리 계산돼요. H1--H3은 대부분 만든 방식 때문에 저절로 성립해요. 그리고 이 가설들은 측정 도구를 점검하는 역할을 해요. 틀렸다고 밝혀질 수 있는 시험은 RQ3--RQ5예요. 그 가운데 마지막 두 질문은, 이 연구가 관련 자료를 뽑기 전에 정한 규칙으로 답해요. 그 규칙은 고친 점수에 다시 적용했어요(\ref{sec:holdout-protocol}절).

\begin{enumerate}
\item[RQ1.] 같은 아비트럼 계약 묶음에서, 허락 화살표만 걷는 엔도스랭크와 송금 화살표만 걷는 AWP는 서로 다른 순서를 만들까요? 그리고 각 점수는 지금의 허락 점검, 그리고 지금의 송금 점검과 얼마나 잘 맞을까요(스피어먼의 로 $\rho$와 켄달의 타우 $\tau$로 재요)?
\item[RQ2.] 같은 계약 묶음에서 각 점수를 계산하는 데, 시간, 메모리, 반복 횟수로 재면 얼마나 들까요? 그리고 두 화살표 집합을 모두 걸으면 얼마나 더 들까요?
\item[RQ3.] 시점 $t_1$에 계산한 점수가 있을 때, 그 점수는 허락을 받는 스펜더에게 들어오는 $t_2$의 새 주인--스펜더 허락과 새 송금자를 예측할까요? 그리고 결합 페이지랭크를 포함해, $t_1$에서 같은 화살표 종류의 원시 차수보다 그것들을 더 잘 예측하는 페이지랭크가 있을까요?
\item[RQ4.] 허락 층과 송금 층 모두 위에서 도는 결합 페이지랭크는, 표본 안(점수를 계산한 자료)에서뿐 아니라 표본 밖(그 뒤의 새 자료)에서도, 어느 한 층만 걷는 걷기보다 계약에 순위를 더 잘 매길까요?
\item[RQ5.] AWP처럼 송금을 걷는 걷기에 허락 층을 더하면, 표본 밖에서 그 걷기가 나아질까요?
\end{enumerate}

RQ3은 같은 두 구성개념의 나중 라벨을 써요. 여기서 순위를 매기는 계약들에는 빚을 갚지 못했는지(채무 불이행)나 빚을 갚았는지(상환)를 나타내는 라벨이 없기 때문이에요(\ref{sub:non-claims}절). RQ1과 RQ2는 두 한 층 걷기와 그 비용을 밝혀요. RQ4와 RQ5는 두 화살표 종류를 견주는 데서 나아가, 둘을 합치는 쪽으로 넘어가요. RQ5의 규칙인 규칙 B는 W0 라벨을 본 뒤에 정했어요. 그래서 규칙 B는 W1에서만 판정해요(\ref{sec:holdout-protocol}절과 \ref{sub:fresh-holdout}절).

\dissertationSubheading[sec:literature-summary]{문헌 검토 요약}

\ref{ch:literature}장은 네 분야를 살펴봐요. 그래프 순위 매기기에서, 페이지랭크 \parencite{page1999pagerank}는 되풀이해서 퍼뜨린 보증으로 점에 순위를 매겨요. 이때 반복 한 번에 드는 비용은 화살표 수에 비례해서 늘어나요 \parencite{langville2006pagerank}. 아이겐트러스트 \parencite{kamvar2003eigentrust}와 트러스트랭크 \parencite{gyongyi2004trustrank}는 정규화(합이 1이 되게 맞추기)와 씨앗 집합을 써서, 페이지랭크를 피어 투 피어 망의 신뢰와 웹 스팸에 맞게 바꾸었어요. 온체인 평판에서, \textcite{do2023hodgerank}의 연구는 이더리움 거래에 페이지랭크와 호지랭크를 적용했어요. AWP \parencite{do2023awp}는 거래 화살표에 금액과 나이에 따라 무게를 달아요. 그리고 자기 순위를 들어오는 송금의 개수와 금액과 견주어 점검해요. 이 통계는 AWP의 그래프와 같은 자료에서 나와요. \textcite{lin2025riskprop}에서는 계정의 위험을 거래 네트워크를 따라 퍼뜨려요. 블록체인 믿음·평판 시스템을 살펴본 조사 \parencite{bellini2020trust}와 스마트 계약 위에 만든 평판 시스템을 살펴본 조사 \parencite{almasoud2020reputation}는, 송금이나 직접 매긴 평가에 기대는 설계들을 설명해요. 이와 달리, 휴먼 패스포트(Human Passport) 같은 이미 쓰이는 디파이 도구는 온체인과 오프체인의 자격 증명을 모아서, 한 주소가 서로 다른 한 사람의 것인지 정해요 \parencite{humanpassportnddocs}. 이 연구들 가운데 어느 것도 ERC-20 허락 \parencite{vogelsteller2015erc20}을 순위 매기기 화살표로 쓰지 않아요. 주소 묶기(같은 주체의 주소를 한데 모으기)는 허락을 한 주체의 주소들을 잇는 데에만 써요 \parencite{victor2020clustering}. 디파이의 위험과 측정에서, 지식을 체계적으로 정리한 연구들은 위험을 서비스의 층에 따라 \parencite{werner2022sok}, 그리고 실제로 일어난 공격에 따라 \parencite{zhou2023sok} 나누어요. 자전 거래(자기끼리 사고팔아 거래가 많은 척하기) \parencite{victor2021washtrading,cong2023washtrading}, 추출 가능한 가치(거래 순서를 이용해 빼내는 이익) \parencite{qin2022bev}, 사기 토큰 \parencite{xia2021scamtokens}, 에어드롭 파밍(공짜 토큰을 노린 가짜 활동) \parencite{luo2025airdrop}에 대한 연구는 송금 활동의 많은 부분이 꾸며 낸 것임을 보여 줘요. 측정 이론에서, 구성 타당도 \parencite{cronbach1955construct,messick1995validity}는 점수가 재는 것과 바깥 기준을 구분해요. 그리고 부트스트랩 \parencite{efron1979bootstrap,diciccio1996bootstrap}은 순위 통계에 대해, 분포를 미리 가정하지 않는 구간을 줘요. 비어 있는 부분(\ref{sec:research-gap}절)은, 허락 화살표와 송금 화살표를 결합하는 걷기를 한 코호트에서 각 화살표 종류를 걷는 걷기 및 원시 차수와 나란히 평가하는 일이에요. 이 평가에는 신뢰구간(믿을 만한 범위), 짝지은 차이, 미리 정한 시간 재기 절차, 시간 홀드아웃이 들어가요.

\dissertationSubheading[sec:methodology-summary]{연구 방법 요약}

\ref{ch:methodology}장은 공개된 장부 자료에, 수로 재고 서로 견주는 연구 설계를 써요. 2023년 10월 1일부터 2026년 9월 30일까지 아비트럼 원의 ERC-20 \texttt{Approval} 로그와 \texttt{Transfer} 로그는 구글 빅쿼리(Google BigQuery)에서 두 번에 나누어 훑어서 뽑아요. 두 번째는 W1을 등록한 뒤에 뽑았어요. 주된 분석은 각 동결일 전 열두 달 동안의 다섯 토큰 이벤트 가운데, 그 동결일의 코호트 둘레에 있는 것만 남기고, 그 금액을 달러로 매겨요($T_{\mathrm{obs}}$에서는 그 코호트가 매칭 코호트이고, $n=6{,}483$이에요). 이 선택들은 등록 분석의 결과를 안 뒤에 했어요(\ref{sec:amount-revision}절). 거듭제곱 반복(같은 계산을 되풀이하기)으로 푸는 페이지랭크 구현 하나가, 두 층을 모두 걷는 결합 걷기, 마지막 허락 그래프, 송금 그래프에 점수를 매겨요. 각 점수는 스피어먼의 로($\rho$)와 켄달의 타우($\tau$)로 견주어요. 같은 기간에서는 송금, 허락, 시빌 안정성 대리 지표와 견주어요. 기간 밖에서는 W0과 W1의 스펜더 계약에서 새 허락 쌍, 새 송금자와 견주어요. 이때 각 동결일의 원시 차수를 기준선으로 써요. 금액 기울기 점검과 표본 정의 점검을 빼면, 모든 계수에는 계약을 짝지어 다시 뽑은 표본 400개로 구한 95\% 구간이 있어요(시드 42). 방법끼리의 비교는 짝지어서 해요. W1과 함께 등록한 비교 하나는, GMX~V2 계약 트레이더가 나중에 청산을 피하는지를 두고 허락 쪽 점수와 송금 쪽 점수를 두 기간 모두에서 견주어요. 이 비교에서는 짝지어 다시 뽑은 표본 2{,}000개를 써요. 그리고 두 판정 대비에는 본페로니 97.5\% 구간(여러 번 비교할 때 넓힌 신뢰구간)을 써요. 모든 표의 바탕이 되는 코드, 설정, 컴퓨터가 읽을 수 있는 요약은 공개 저장소에 있어요(\ref{sec:reproducibility}절). 그리고 이 연구는 공개된 자료만 써요(\ref{sec:ethics}절).

\dissertationSubheading[sec:contributions]{기여}

이 논문은 온체인 평판과 디파이 위험 연구에 다섯 가지 기여(이 연구가 새로 보탠 것)를 해요. 엔도스랭크와 AWP는 둘 다 무게를 단 페이지랭크를 한 가지 화살표 종류에 적용해요 \parencite{page1999pagerank}. 둘은 각자 자기 재시작 분포(어디서 다시 출발하는지의 비율)를 써요. 결합 걷기인 C-PR은 이동 규칙(다음에 어느 점으로 갈지 정하는 규칙) 그 자체를 바꿔요. 그래서 새로운 것은 다음과 같아요. 두 가지 온체인 관계에 걸친 이동 규칙, 미리 정한 규칙으로 한 그 이동 규칙의 시험, 화살표 종류 하나, 평가 설계, 그리고 이런 걷기가 이 표본에서 무엇을 계산하는지에 대한 수학적 설명이에요.
\begin{itemize}
\item 허락과 지급을 함께 걷는 결합 걷기(C-PR): 달러로 나타낸 허락 층과 달러로 나타낸 송금 층을 걷는 무작위 걷기 하나예요. 허락 층에서는 허락 하나하나에 상한을 둬요. 이 걷기는 각 주소에서, 그 주소가 들어 있는 층들에 대해 두 이동 행을 미리 정해 둔 섞는 비율 $\lambda$로 섞어요. S-PR은 이 걷기의 씨앗을 쓰는 빼 보기 실험이에요. 살펴본 평판 방법 가운데 이 두 관계를 결합하는 것은 없어요(\ref{sec:research-gap}절).
\item 미리 정한 규칙으로 한 그 걷기의 시험: 규칙 A는 분석 계획에서 $\lambda$와 함께 정했어요. 규칙 A는 그 주된 부분에서, C-PR이 새 허락 쌍에 대해 자기 허락 층보다 못하지 않고 새 송금자에 대해 자기 송금 층보다 낫기를 요구했어요. 규칙 B는 W1과 함께 등록했어요. 규칙 B는 C-PR이 자기 송금 층보다 나아지기를 요구했어요. 결합만 바꾸고 다른 것은 하나도 바꾸지 않는 대비는 이것들뿐이에요. 고친 점수에서 재분석으로 보면, 두 규칙 모두 등록 분석에서처럼 실패해요(부록~\ref{sec:map-registered}). 그래도 새 허락 쌍에서 송금 층보다 나은 점은 W1에서 다시 나타나요(주장~C1). 규칙과 별도로, C-PR은 두 기간 모두에서 두 스펜더 라벨을 엔도스랭크보다 더 잘 줄 세워요. 하지만 나중의 송금자는 AWP보다 못하게, 나중의 허락 쌍은 받은 허락 수보다 못하게 줄 세워요. 그리고 나중의 허락에서 AWP보다 앞서는 것은 대부분 AWP가 순서를 매길 수 없는 계약들에서 와요(주장~C2). 빼 보기 실험은 허락 신호가 어떻게 되는지 보여 줘요. 허락 화살표가 걷기를 직접 이끌 때는 허락 신호가 살아남아요. 허락 화살표가 걷기가 어디서 다시 출발할지만 정할 때는 그 신호가 대부분 사라져요.
\item 마지막 허락으로 화살표를 만드는 방법(엔도스랭크), 곧 허락 쪽 비교 대상: 보증 그래프에서 화살표 $u \to v$는 스펜더 $v$가 주인 $u$에게서 받아 가진, 기간 안의 마지막 ERC-20 허락을 나타내요. 이 허락은 \texttt{approve}나 \texttt{permit}으로 정해져요. 금액은 달러로 매기고, AWP가 송금에 무게를 다는 방식대로 무게를 달아요. 이 화살표는 돈을 낸 기록이 아니라, 쓸 권리를 맡긴 일을 기록해요. 허락은 한 주체의 주소들을 묶는 데 쓰인 적이 있어요 \parencite{victor2020clustering}. 하지만 살펴본 평판 방법 가운데 허락으로 순위를 매기는 것은 없어요(\ref{sec:research-gap}절).
\item 기간 밖 시험과 차수 기준선을 갖춘, 같은 코호트에서의 평가: 모든 점수는 한 계약 묶음에서 계산해요. 그리고 AWP와 함께 쓰인 절차에 따라, 똑같은 송금, 허락, 시빌 안정성 대리 지표로 점검해요 \parencite{do2023awp}. 계수에는 95\% 부트스트랩 구간을, 대비에는 짝지은 구간을 구해요 \parencite{efron1979bootstrap}. $t_1$에 동결한 점수는 $t_1$ 뒤에 생기는 새 허락과 새 송금자로 시험해요. 이때 페이지랭크가 다듬는 원시 차수도 나란히 시험해요. 이렇게 해서 표본 안의 일치와 예측을 나누고, 퍼뜨린 점수와 그 점에서 그냥 센 수를 나누어요(주장~C2와~C3). 라벨을 본 뒤에 계산한 대비에서, 기간 밖에서 자기 차수를 이긴 페이지랭크는 없었어요. AWP는 나중의 송금자를 송금 입차수보다 못하게 줄 세웠어요. 엔도스랭크는 어느 다시 출발 규칙을 쓰든 나중의 허락 쌍을 받은 허락 수보다 못하게 줄 세웠어요. C-PR은 어느 라벨도 같은 화살표 종류의 차수보다 잘 줄 세우지 못했어요. W1과 함께 등록한 비교 하나는 어느 화살표 종류로도 만들지 않은 결과를 더해요. 나중에 청산을 피하는 것에서, 등록 기간에는 허락한 주인 수가 송금 수보다 앞섰어요. 하지만 W0에서는 어느 쪽이 나은지 가려지지 않았어요. 서로 맞대어 볼 수 있게 맞춘 페이지랭크 짝도 가려지지 않아요. 트레이더 표본은 작아요(주장~C6).
\item 깊이 1인 그래프에서의 닫힌 꼴(한 줄 식으로 바로 쓴 답): 고르게 다시 출발하고, 매달린 점(나가는 화살표가 없는 점)의 몫을 고르게 다시 나누면, 들어오는 이웃(화살표를 보내오는 점)들에게 들어오는 화살표가 하나도 없는 스펜더는 점수로 $c\,(1+d\,\delta(v))$를 받아요. 여기서 $c$는 들어오는 화살표가 없는 점들의 공통 점수예요. $d$는 감쇠 계수예요. $\delta(v)$는 $v$로 들어가는 이동 확률의 합이에요. 엔도스랭크에서 이 값은 주인 기준으로 나눈 들어오는 무게예요. 이 값은 각 주인의 보증을 그 주인이 허락한 모든 스펜더에게 나누어요. 이 명제에는 고르게 다시 출발하는 것만 필요해요. 그래서 이 명제는 C-PR의 걷기도 설명해요. 엔도스랭크의 그래프에서, 이 식은 깊이 1인 모든 스펜더에서 1\% 안으로 맞아요. 하지만 더 깊은 스펜더 가운데에서는 소수에서만 맞아요. 그래서 여러 걸음에 걸친 퍼뜨림이 있어요(\ref{sub:theory-closed-form}절). 시간 무게와 토큰 무게 말고도, 그 나누기와 이 퍼뜨림이 엔도스랭크를 그냥 센 허락한 주인 수와 다르게 만들어요. 하지만 전체 스펜더에서는, 그 때문에 엔도스랭크가 나중의 허락을 그 수보다 더 잘 줄 세우게 되지는 않았어요. 그리고 더 깊은 스펜더만 따로 떼어 시험하지는 않았어요.
\end{itemize}


\dissertationSubheading[sec:scope-limitations]{범위와 한계}

이 연구는 2023-10-01부터 2026-06-30까지의 아비트럼 원 자료를 뽑아요. W1은 2026년 7월부터 9월까지를 더해요. 그리고 주된 분석은 각 동결일 전 열두 달로 점수를 매겨요. 주된 표본은 계약 $n=6{,}483$개로 된 매칭 코호트예요. 홀드아웃 표본은 각 동결일의 스펜더 계약이에요(W0에서 $n=7{,}914$, W1에서 $n=7{,}960$). 점수를 받는 지갑은 계약 계정이에요. 사용자 지갑은 쓰임새가 너무 여러 가지라서, 그 활동을 한 줄로 세울 수 없기 때문이에요(\ref{sec:data-sources}절). 계정 종류는 이벤트가 일어난 때마다가 아니라, 정해 둔 블록 하나에서의 코드로 읽어요. 다섯 토큰의 이벤트 가운데 한쪽 끝에 코호트 계약이 있는 것만 그래프에 들어가요. 그래서 각 그래프는 코호트의 에고 네트워크(코호트 계약에 닿는 화살표만 모은 그래프)예요. 그리고 여러 걸음으로 이어진 길은, 그 길의 화살표 하나하나가 코호트 계약에 닿을 때만 보여요(\ref{sub:theory-closed-form}절). 동결일 그래프에 있는 GMX~V2 계약 계정 가운데 이 기간들에 포지션을 세 번 이상 닫은 것은 얼마 없어요(W0에서 $n=87$, W1에서 $n=70$). 그래서 트레이더 결과는 모두 구간이 넓어요(\ref{sub:neutral-label-results}절). 실행 시간은 코호트의 부분 표본에서도 재요. 결과가 아비트럼 원이나 살펴본 몇 달을 넘어서도 들어맞는다고 주장하지 않아요. 달러로 나타낸 금액과 열두 달 점수 기간은 등록 분석의 결과를 안 뒤에 골랐어요(\ref{sec:amount-revision}절). 엔도스랭크와 AWP는 금액을 상한이 있는 함수에 넣어요. 이 함수는 중앙값 금액에 절반의 무게를 줘요. 반면 C-PR의 층들은 금액 그 자체를 쓰고, 허락 하나하나에는 상한을 둬요(\ref{ch:methodology}장). 허락이나 송금을 나쁜 뜻으로 조작하는 것은 실제로 있는 위협이에요. 주소를 만드는 데 돈이 거의 들지 않기 때문이에요 \parencite{douceur2002sybil}. \ref{sec:sybil-cost-model}절은 다음을 보여 줘요. 작은 시빌 주소 농장(공격자가 만든 가짜 주소 무리)은 적은 비용으로 허락 그래프 위에서 한 계약의 점수를 코호트 대부분보다 위로 올릴 수 있지만, 꼭대기에 오르는 것은 비싸요. 이것을 뒷받침하는 한계 식(점수가 얼마까지 오를 수 있는지 정하는 식)은 C-PR에도 들어맞아요. 그리고 시빌 보정 대리 지표는 몇몇 조작 방식을 보여 주지만, 점수를 그런 조작에서 지켜 주지는 않아요.

\dissertationSubheading[sec:roadmap]{연구 길잡이}

\ref{ch:literature}장은 그래프에 바탕을 둔 온체인 평판, ERC-20 허락의 뜻, 디파이에서 쓰이고 있는 평판 도구를 살펴봐요. 그리고 이론의 틀을 세워요. \ref{ch:methodology}장은 코호트, 토큰, 가격과 금액의 수정, 그래프와 점수, 구성개념 점검, 시간 홀드아웃, 시간 재기 절차, 그리고 재현성(다른 사람이 똑같이 다시 해 볼 수 있음)과 윤리의 기준을 정의해요. \ref{ch:results}장은 다음을 보고해요. 벤치마크(속도 재기, 주장~C5), 모든 점수의 같은 기간 일치도와 두 한 층 걷기의 순위 갈라짐(C4), 강건성 점검(조건을 바꿔도 결과가 버티는지 보기), 한 층 걷기들의 시간 홀드아웃(C3), 규칙에 따라 본 결합 걷기와 다른 모든 점수 옆에 놓고 본 결합 걷기(C1과 C2), 그리고 청산을 피하는 것에서의 비교(C6)예요. 부록~\ref{ch:appendix-er-awp}에는 엔도스랭크와 AWP를 직접 견준 결과를 모았어요. 부록~\ref{sec:map-registered}은 등록 분석을 보고해요. \ref{ch:discussion}장은 찾아낸 것을 연구 질문과 목표에 따라 풀이해요. 그 뜻하는 바를 따져 보고, 허락 그래프에 대한 시빌 공격(가짜 주소를 많이 만들어 속이는 공격)의 비용을 모형으로 만들어요. 그리고 타당성을 위협하는 것과 한계를 따져 봐요. \ref{ch:conclusion}장은 RQ1--RQ5에 답하고, 지식에 대한 기여를 밝히고, 앞으로 할 일을 내놓아요.
